\section{Introduction}\label{sec:intro}

Position, navigation, and timing (PNT) underpins power-grid synchronisation,
financial timestamping, telecommunications, and essentially every form of
autonomous and crewed mobility. In practice that infrastructure leans almost
entirely on global navigation satellite systems (GNSS), whose signals arrive at
the Earth's surface tens of decibels below the thermal noise floor and are
therefore trivially denied: low-cost jammers, meaconing, and increasingly
sophisticated spoofers can suppress or counterfeit a GNSS fix over a wide
area~\citep{misra2010,kaplan2017}. The operational consequence is loss of a
\emph{reference}, not merely of a gadget. The question that matters during an
outage is how long a platform can hold its position and time to within a stated
tolerance once the satellites can no longer be trusted. We refer to this
property as \emph{PNT resilience}: the growth of position and timing error
against outage (or holdover) duration, for a given on-board sensor suite.

A platform rides through a GNSS outage on whatever it carries. Inertial sensors
let it dead-reckon, but their error grows super-linearly as bias and
random-walk noise integrate twice into position~\citep{groves2013}; a local
clock holds time, but only until its own instability accumulates past the
allowed offset~\citep{riley2008nist}. The classical floor for both is well
characterised. The open question, and the subject of sustained investment, is
whether \emph{quantum sensing} meaningfully raises that floor. Cold-atom
interferometric accelerometers and gyroscopes promise bias stability and scale
factors set by atomic constants rather than by a drifting proof mass; optical
atomic clocks reach fractional frequency instabilities orders of magnitude below
chip-scale or even flight-qualified microwave standards; and optical two-way
time transfer offers a path to disseminate that stability between platforms.
The common thread is a \emph{slower error-growth law}: if the per-second noise
and the long-term drift are smaller, the error crosses any fixed tolerance later,
extending the survivable outage. Recent reviews survey this case in
detail~\citep{coldatomnavreview2022,quantumpntreview2026}. Tellingly, those same
reviews concede that quantitative comparisons across the field are difficult to
audit: they frequently rest on generic synthetic scenarios, mismatched figures
of merit, and proprietary trade studies whose assumptions are not published.
That concession matters, because it names the gap precisely: the resilience
debate is starved not of physics but of \emph{reproducible, falsifiable,
like-for-like comparison}.

The gap is real, and it is worth establishing by evidence. A number of capable
open tools exist, but each answers a different question and
none reproduces an end-to-end quantum-versus-classical resilience comparison on a
single neutral code path. RTKLIB~\citep{takasu2009} is a mature open package for
processing real GNSS observations into precise positions, but it concerns the
fix itself, not the post-outage error-growth dynamics of an inertial or clock
holdover. \texttt{gnss\_lib\_py}~\citep{knowles2024gnsslibpy} provides a clean
Python substrate for GNSS measurement analysis and algorithm prototyping, again
on the receiver-processing side rather than the resilience side. LuPNT
\citep{iiyama2023lupnt} targets lunar PNT architecture studies and is the
closest in scope to the multi-domain resilience question, yet it is an
architecture simulator, not an instrument for swapping a quantum sensor model
against its classical counterpart under controlled uncertainty. NaveGo
\citep{gonzalez2017navego} is a well-validated open inertial-navigation toolkit
with published error profiles (genuinely strong on the inertial axis), but it
models classical IMUs and does not carry clock-holdover, time-transfer, or
quantum-sensor models. None of these is unfalsifiable; the point is the
converse. Each is a sound tool for its own question, and the absence is the
\emph{intersection}. To our knowledge no open tool runs the full PNT resilience
stack of inertial dead-reckoning, clock holdover, and time transfer with a
quantum and a classical configuration exercised through the same code path,
seed-reproducibly, with uncertainty propagated.

This paper presents Kshana, an open engine built to close that intersection, and
uses it to produce a genuinely new result. The engine's external-validation
battery (Section~\ref{sec:validation}) establishes that its underlying
mechanics are trustworthy: frames bit-for-bit against SOFA/ERFA, SGP4 to
millimetre agreement against the AIAA reference set, Allan deviation against
NIST and Stable32, and force-model validation by ephemeris fitting against
agency precise orbits. That battery is \emph{evidence}, not the
contribution. The contribution is the reproducible crossover study built on top
of it (Section~\ref{sec:crossover}). Instead of reporting a single hero number,
we sweep each quantum lever across its published-parameter envelope by
Monte-Carlo and report the resilience advantage as a distribution with 95\%
confidence intervals, then locate the \emph{break-even} against the classical
baseline. The result is a map of \emph{when} quantum sensing changes the outcome
and when it does not. For the inertial lever, the cold-atom accelerometer's
advantage over a navigation-grade IMU is decisive at low platform vibration but
the break-even migrates toward higher vibration as the outage lengthens, so the
honest answer is conditional on the host platform's mechanical environment. For
the clock lever, an optical-lattice standard and a flight-qualified maser hold
sub-microsecond time across a full 24~h holdover in our envelope, whereas a
deployed chip-scale clock crosses the same threshold near eight and a half
hours: a same-maturity comparison that separates what is demonstrated in the
lab from what flies today. We label every quantum result as a performance-model
projection, not a measurement, and carry its technology-readiness tag inline.

We claim no superiority on any single axis, and we use ``to our knowledge, no
open tool \ldots'' rather than ``first'' or ``only.'' In-browser orbit
propagation, for instance, is not new~\citep{satellitejs}; our delta is that a
single engine runs the full integrated stack zero-install. With that framing,
our contributions are:

\begin{enumerate}
  \item \textbf{A reproducible quantum-versus-classical resilience crossover map
  under parameter uncertainty} (the new result, Section~\ref{sec:crossover}).
  For each quantum lever (cold-atom inertial, optical-lattice clock, optical
  time transfer), a Monte-Carlo sweep over the published-parameter envelope
  yields the resilience-figure-of-merit advantage as a distribution with 95\%
  confidence intervals and identifies the break-even against the classical
  baseline, answering \emph{when} quantum sensing changes the outage outcome,
  openly and from a committed scenario and seed.

  \item \textbf{A reproducibility-and-falsifiability methodology} for
  resilient-PNT comparison (Section~\ref{sec:method}): a scenario-plus-seed-plus-
  version reproducibility triple, a per-figure-of-merit \emph{validated} or
  \emph{not-modeled} label, and hand-derived oracle tests. It builds on prior
  reproducible-research practice~\citep{fernandezprades2018reproducibility,nasastd7009};
  the genuinely new element is the requirement that the quantum and classical
  configurations differ only by a parameter swap on a \emph{single neutral code
  path}, so that any reported advantage is attributable to the device model and
  not to two divergent implementations.

  \item \textbf{One open engine instantiating the methodology across the PNT
  stack} (Section~\ref{sec:arch}): a Rust core with native, Python, and
  WebAssembly targets, 19 scenario kinds, and standard-format interoperability
  (RINEX, SP3, CCSDS OEM/OMM), summarised with its capabilities \emph{and}
  non-capabilities stated symmetrically in Table~\ref{tab:intersection}.

  \item \textbf{An external-validation battery, tiered by strength}
  (Section~\ref{sec:validation}, Table~\ref{tab:validation}): bit-for-bit checks
  against authoritative oracles (frames, SGP4, Allan deviation); post-fit
  residuals against real agency precise ephemerides, with the force-model-fit,
  arc-length, and gravity-degree caveats stated inline (full-arc Galileo
  $0.61$~m at degree/order~70 as the headline; the reduced-dynamic
  empirical-acceleration tier reported as a labelled bound, not a measurement of
  force-model accuracy; LRO $6.6$~m, above the $5$~m target, stated honestly);
  and a controlled, reproducible cross-validation showing the kernel-free
  analytic stack reproduces a DE440/ANISE realisation for the LRO reduced-dynamic
  arc, consistent with established results~\citep{bertone2021grail,li2023leopod}
  and presented as a reproduction scoped to that arc, not a discovery.

  \item \textbf{A multi-layer spoofing detector}, characterised against published
  TEXBAT \emph{parameters} (Section~\ref{sec:validation}), with raw-IQ ingestion
  stated as future work; it sits explicitly in the not-externally-validated tier.
\end{enumerate}

The remainder of the paper develops the methodology that makes these claims
falsifiable (Section~\ref{sec:method}), situates the engine against related open
tools (Section~\ref{sec:related}), summarises its architecture
(Section~\ref{sec:arch}), presents the crossover study
(Section~\ref{sec:crossover}) and the tiered validation
(Section~\ref{sec:validation}), and states the limitations and threats to
validity (Section~\ref{sec:limitations}), chief among them that the
quantum-sensor figures throughout are model projections from cited device
literature, not measurements.
